\section{Weekly Summary}

\begin{tcolorbox}[colback=yellow!10,colframe=black!50,title=AI-Generated Content]
\noindent The summary below was written by Claude (Anthropic) from that week's regression outputs and series descriptions. It has not been reviewed or fact-checked by a human, and should be read as an automated, informal recap -- not analysis.
\end{tcolorbox}

Welcome back to the weekly parade of economic data being forced into relationships it never consented to. This week the random number generator rolled up five pairings, and as always, we ran each one twice: once raw, and once after stripping out the shared march of time that makes everything look suspiciously friendly. Grab a beverage. It's a journey.

We kicked things off Monday with the all-time banger matchup of Light Truck Production Seasonal Factors versus the Export Price Index for Nonferrous Metals That Aren't Aluminum. The suspense was unbearable, and then the result arrived: an R-squared of basically zero, a p-value of 0.93, and a coefficient so limp it might as well have phoned in sick. Both raw and detrended agreed on absolutely nothing happening. A perfect nothingburger, cooked to order. Tuesday wasn't much better -- Latin American ag import prices versus Checkable Deposits Held By The 50th-to-90th Wealth Percentiles (a series so discontinued it's practically a ghost). The raw regression lit up with a thrilling p-value of 0.0006 and I could feel myself reaching for the champagne, but then detrending walked in, yanked out the shared upward drift, and the whole thing collapsed into a p-value of 0.86. Classic trend mirage. Two things going up over time is not a love story, it's a coincidence with good posture.

Wednesday's pairing is the one that'll haunt my dreams: the DC minimum wage versus the Peoria, Illinois house price index, raw R-squared of 0.77 and a p-value with so many zeros it looked like a typo. For one shining moment I thought we'd discovered that DC wage law secretly governs Peoria real estate. Then detrending dropped it to a p-value of 0.91 and an R-squared of 0.0003, i.e., nothing, because of course both things just grew over the decades like everything else priced in dollars. Thursday actually gave us the week's rarest creature: Midwest regular gasoline prices versus Japanese industrial producer prices held up in BOTH the raw (p around 0.0002) and the detrended version (p around 0.025). It's weak -- the detrended R-squared is a modest 0.05 -- but it survived the one test that usually kills these things, and energy prices genuinely do ripple through industrial production costs worldwide, so this one gets a grudging raised eyebrow rather than an eye-roll. Friday closed us out with the Smoothed World Uncertainty Index for Yemen versus discontinued Plastics import prices, which flatlined in both versions and went quietly into the night.

As always, the mandatory disclaimer that I will keep repeating until I'm blue: these are randomly paired series thrown at a regression for entertainment. Nothing here is a causal claim, a forecast, or anything you should mention in a serious conversation. Spurious correlation isn't the risk in this project -- it's the entire business model. The gasoline-and-Japan result is mildly interesting precisely because it's the exception, and even that I'd trust about as far as I could throw a discontinued FRED series. See you next week for more statistics behaving badly.
